\scanpage{231}\mspage{86}
\deleted{À dire vrai} Notons que si \uline{$W$ est stable par composition} (ce que
nous allons supposer désormais), il en est de même
de $\mathbb A_W(L)$ (ce qui utilise \deleted{de façon} encore
le fait que $L$ soit stable par changt de base).
De plus, si \uline{$W$ contient les iso} (ce que nous
supposerons également), $\mathbb A_W(L)$ les contient aussi.
Quant à $\mathbb L_W(A)$, il contient les iso.
\uline{pourvu que} $A$ soit stable par isomorphisme\unsure{;}
quand $A$ est stable par composition et contient les
\unsure{iso|id}, cette condition sera automatiquement satisfaite. Nous tenons à cette condition, c'est
pourquoi du côté de $A$, dans la définition
de $P$, \unsure{je} vais introduire cette condition
supplémentaire que $A$ soit \uline{stable par iso}

Donc je pose, en corrigeant (6.2)
\[
\left\{\begin{aligned}
P={}&\{A\in\unsure{\mathfrak{P}|\mathcal{P}}(\operatorname{Fl}(M))\mid
\deleted{A\subset W,\ \unsure{A}\ A\ \text{stable par iso}\}}\\
&\inserted{\operatorname{Is}(M)\subset A\subset W,\ A\ \text{stable par composition}}\},\\
\deleted{\text{\unsure{Notons que}}}\quad
Q={}&\{L\in\mathfrak{P}(\operatorname{Fl}(M))\mid\operatorname{Is}(M)\subset L,\ L\ \text{stable }\deleted{\text{par}}\\
&\qquad\deleted{\text{composition et}}\ \text{par chgt. de base}\}.
\end{aligned}\right.\tag*{(6.7)}
\]

On est bien sûr intéressé surtout aux
\unsure{éléments|élémts} clos de $P,Q$ i.e. les parties
$A,L$ de $\operatorname{Fl}(M)$, \unsure{contnt} les iso, stables
par composition, (et $L$ stable aussi par chgt
de base, $A\subset W$) telles que l'on ait
\[
A=\mathbb A_W(L),\quad L=\mathbb L(A)\tag*{(6.8)}
\]
% archive sheet number: 231
